
Group Relative Policy Optimization (GRPO) has emerged as a method for training code-generating language models from execution feedback, yet its effectiveness depends on a structural precondition that is not always verified in practice: at least one completion in each training group must receive non-zero reward for advantage normalization to carry gradient signal. This paper examines what happens when that precondition is systematically violated. In a controlled comparison of function-level (APPS+CodeContests) and repository-level (SWE-bench Verified) GRPO training under otherwise identical configuration, we observe a 76× advantage variance gap (0.004 vs. 0.317, Welch's t-test: t=20.37, p=5.34×10⁻⁴⁴, Cohen's d=1.904) using CodeLlama-7b-Instruct-hf. This gap is explained by the difference in positive reward rates: approximately 0\% for function-level competitive programming tasks versus approximately 6\% for repository-level tasks. When positive reward rate is near zero, every training group is degenerate (std(r)=0), advantages are zero, and GRPO gradient contributions vanish. We additionally validate training infrastructure for a 4-condition difficulty-stratified curriculum GRPO experiment (curriculum, uniform, easy-only, hard-only) on APPS+CodeContests via a 10-step smoke test; reward density was 0.0 across all conditions and steps during this smoke test, consistent with the reward collapse finding. The full 5000-step behavioral experiment was not executed. Our results characterize advantage variance as a diagnostic for execution-feedback RL and identify conditions under which curriculum ordering may or may not rescue degenerate training.
